\subsection{Ordinary backtracking and an auto-conditioned comparator}
A separate 350-run extension uses the same seven families, ten seeds,
tolerance, and cap. It compares ordinary backtracking initialized at 1
or $U_0/D^2$, certificate scaling, Armijo, and AC-FW
\cite{autoconditioned}. Of these runs, 280 repeat existing deterministic
trajectories with additional logging; only the 70 AC-FW runs add a new
method on these same instances. They are not new independent datasets.

AC-FW is our independent implementation of Algorithms 1 and 2 in
\cite{autoconditioned}, using its equation (6) with a fixed $\delta=1$:
\[
 r_t=1-\frac{1}{(t+1)\log^2(t+3)}.
\]
We do not claim to reproduce the authors' numerical configuration.
Its compulsory move from $x_{-1}$ to an LMO vertex $x_0$, both function
and gradient evaluations, and the initial LMO are charged. Available
linear lower bounds from both points enter the certificate history.
After initialization, each outer iteration makes one function trial;
strict decrease accepts the point, and rejection keeps the old point
and cached gradient. Rejections consume the iteration cap. Consequently
$N_f=N+2$, $N_{\rm LMO}=N+2$, and $N_g=2+N_{\rm accepted}$, including the
final certificate check. The initial move changes the starting trajectory;
its effect is included in the algorithm comparison, not isolated in an
initialization ablation. Our cosh and Poisson losses are not globally
gradient-Lipschitz on $\R^d$, unlike the assumption in the cited paper;
their use here is an empirical comparison on the compact feasible set,
not a claim that its global theorem applies verbatim.

Table~\ref{tab:backtracking} shows strong competing results. On the active
family AC-FW certifies all ten runs, with medians 18.5 function and gradient
calls; certificate scaling certifies one run within the cap. On the scaled
active family neither adaptive variant nor AC-FW certifies any run, and
AC-FW's median certificate is worse. On Bike Sharing AC-FW uses 639 median
function calls and certifies seven starts, versus 860 calls and eight
starts for certificate scaling; its median gradient count is higher
(639 versus 420). Capped runs remain in every median, so these figures
must be read together with successes and both oracle budgets. Only five
outer trials are rejected across the 70 AC-FW runs. This study gives no
uniform practical ranking or evidence that the normalization is necessary.

We replay all 140 ordinary-backtracking runs from their archived source
and check every accepted curvature, decrease, phase classification,
lag budget, telescoping trial identity, and trial bound. The replay
matches recorded objectives, certificates, and all three oracle counts
exactly. All these runs have zero lag steps relative to the conservative
analytic phase caps; they do not test sharpness of the lag bound or the
large worst-case iteration constants. A separate affine example tests
more than twenty lag steps and the exact formula in
Proposition~\ref{prop:initialization}. Five hundred abstract trajectories
stress the potential argument with abrupt phase jumps and varying
initial scales. These checks support implementation consistency, not
independent mathematical validation or a priority claim.

\begin{table}[htbp]
\centering\small
\caption{350-run extension, selected families. Medians include capped runs. The cap is 1000 outer iterations, including AC-FW rejections; the compulsory AC-FW initialization is included in all oracle costs. Success means $U\le10^{-5}$.}
\label{tab:backtracking}
\begin{tabular}{lrrrrr}
\hline
Method & Successes & Outer iterations & $U$ & $f$ calls & $\nabla f$ calls \\
\hline
\multicolumn{6}{l}{\textbf{active}} \\
Ordinary, $M=U_0/D^2$ & 1/10 & 1000.0 & 2.43e-04 & 2002.0 & 1001.0 \\
Ordinary, $M=1$ & 0/10 & 1000.0 & 2.23e-04 & 1998.5 & 1001.0 \\
Certificate scaling & 1/10 & 1000.0 & 1.99e-04 & 2014.0 & 1001.0 \\
AC-FW, $\delta=1$ & 10/10 & 16.5 & 7.93e-06 & 18.5 & 18.5 \\
Armijo & 10/10 & 19.0 & 6.55e-06 & 141.0 & 20.0 \\
\hline
\multicolumn{6}{l}{\textbf{active\_scaled}} \\
Ordinary, $M=U_0/D^2$ & 0/10 & 1000.0 & 1.87e-03 & 2002.0 & 1001.0 \\
Ordinary, $M=1$ & 0/10 & 1000.0 & 1.77e-03 & 2002.0 & 1001.0 \\
Certificate scaling & 0/10 & 1000.0 & 1.47e-03 & 2014.0 & 1001.0 \\
AC-FW, $\delta=1$ & 0/10 & 1000.0 & 2.67e-03 & 1002.0 & 1002.0 \\
Armijo & 1/10 & 1000.0 & 4.78e-04 & 9758.0 & 1001.0 \\
\hline
\multicolumn{6}{l}{\textbf{bike\_poisson}} \\
Ordinary, $M=U_0/D^2$ & 6/10 & 500.0 & 9.87e-06 & 1002.5 & 501.0 \\
Ordinary, $M=1$ & 8/10 & 893.5 & 9.85e-06 & 1789.0 & 894.5 \\
Certificate scaling & 8/10 & 419.0 & 9.41e-06 & 860.0 & 420.0 \\
AC-FW, $\delta=1$ & 7/10 & 637.0 & 9.53e-06 & 639.0 & 639.0 \\
Armijo & 10/10 & 310.0 & 9.93e-06 & 4399.0 & 311.0 \\
\hline
\end{tabular}
\end{table}

